\documentclass[11pt,a4paper]{article}

\usepackage[margin=2.5cm]{geometry}
\usepackage{amsmath,amssymb,amsfonts}
\usepackage{algorithm}
\usepackage{algpseudocode}
\usepackage{graphicx}
\usepackage{bm}
\usepackage{cite}
\usepackage{hyperref}
\usepackage{booktabs}
\usepackage{multirow}
\usepackage{array}

\newcommand{\R}{\mathbb{R}}
\newcommand{\C}{\mathbb{C}}
\newcommand{\erfi}{\mathrm{erfi}}
\newcommand{\erf}{\mathrm{erf}}
\newcommand{\TV}{\mathrm{TV}}
\newcommand{\diag}{\mathrm{diag}}
\newcommand{\Tr}{\mathrm{Tr}}
\newcommand{\sgn}{\mathrm{sgn}}
\newcommand{\sinc}{\mathrm{sinc}}
\DeclareMathOperator*{\argmin}{arg\,min}

\title{Distributed Phased Array Forward Scatter Radar:\\
Mathematical Model and Signal Processing Algorithms\\[6pt]
\large Version 6 --- 2D Receive Arrays for Altitude Estimation}

\author{D.~Huang}
\date{\today}

\begin{document}
\maketitle

\begin{abstract}
This document extends the v5.2 mathematical model by replacing each
one-dimensional ($N$-element) receive array with a two-dimensional
$N\times N$ planar array in the $x$--$z$ plane.  The 2D array
provides both azimuth (direction cosine~$\ell_x$) and elevation
(direction cosine~$\ell_z$) measurements of the target.  In
\textbf{Zone~B}, the elevation DOA yields a direct, coarse estimate
of the target altitude~$z_T$, removing the reliance on amplitude-based
methods or the default $z_T{=}0$ initialisation.  In \textbf{Zone~A},
the improved $z_T$ initial estimate accelerates convergence of the
joint gradient-descent retrieval.  The 2D array also increases
beamforming gain by a factor of~$N$ (from $N$ to $N^2$ elements),
improving SNR in both zones.

All other components --- 3D target model with oblique motion,
pixel-dependent effective propagation distance, phase-compensated
covariance for time-varying azimuth DOA, gradient descent with TV
denoising, and joint parameter estimation --- are retained from v5.2.
\end{abstract}

\tableofcontents
\newpage

%======================================================================
\section{System Geometry}
\label{sec:geometry}
%======================================================================

\subsection{Coordinate System}

Cartesian system $(x,y,z)$: incident plane wave travels in~$+y$.
A set of $K$ phased array receiver nodes is positioned along the
$x$-axis at $y=R_0$.  Node~$k$ is centred at
$(\bar{X}_k,\,R_0,\,0)$, $k=0,\ldots,K{-}1$.

\subsection{2D Planar Array}
\label{sec:2d_array}

Each node has $N\times N$ elements arranged in a uniform rectangular
array (URA) in the $x$--$z$ plane, with inter-element spacing
$d=\lambda/2$ in both dimensions.  Element $(n_x,n_z)$ of node~$k$ is
at
\begin{equation}
  \mathbf{p}_{k,n_x,n_z}
  = \bigl(\bar{X}_k + n_x\,d,\;R_0,\;n_z\,d\bigr),
  \quad n_x=0,\ldots,N{-}1,\quad n_z=0,\ldots,N{-}1.
  \label{eq:element_pos}
\end{equation}
All distances are normalised to $\lambda$, so $\lambda=1$ and
$d=\tfrac12$.  Each node has $N^2$ elements in total.

\paragraph{Comparison with v5.2.}
In v5.2, each node had $N$ elements along~$x$ only (a ULA), providing
azimuth (direction cosine~$\ell_x$) but no elevation information.
The 2D array adds $N$ rows in the $z$-direction, enabling
measurement of the elevation direction cosine~$\ell_z$, which is
directly related to the target altitude~$z_T$.

\subsection{Node Placement}

The $K$ nodes are uniformly spaced at half the Fresnel zone radius:
\begin{equation}
  \Delta X = \frac{\sqrt{R_0}}{2},
  \qquad
  \bar{X}_k = \Bigl(k - \frac{K{-}1}{2}\Bigr)\frac{\sqrt{R_0}}{2},
  \quad k=0,\ldots,K{-}1.
  \label{eq:node_positions}
\end{equation}

\subsection{Target Model --- Oblique Motion}
\label{sec:target_model}

(Unchanged from v5.2.)
The target is a \textbf{three-dimensional rigid body}.  Its centre
moves along a straight line in the $x$--$y$ plane at speed~$v$ and
heading angle~$\theta$:
\begin{equation}
  \text{Target centre at time } t:\quad
  \bigl(\,x_T(t),\;y_T(t),\;z_T\,\bigr),
  \label{eq:target_pos}
\end{equation}
where
\begin{align}
  x_T(t) &= x_0 + v_x\,t, &\quad v_x &= v\cos\theta,
  \label{eq:xT}\\
  y_T(t) &= v_y\,t,        &\quad v_y &= v\sin\theta,
  \label{eq:yT}
\end{align}
with $x_0$ the initial $x$-coordinate, $y_T(0)=0$, and $z_T$ a
constant $z$-offset (the target altitude).

\subsection{Time-Varying Propagation Distance}
\label{sec:R_of_t}

(Unchanged from v5.2.)
\begin{equation}
  R(t) = R_0 - v_y\,t = R_0 - v\sin\theta\;\!t.
  \label{eq:R_of_t}
\end{equation}

%======================================================================
\section{3D Target Shape and Shadow Projection}
\label{sec:3d_shadow}
%======================================================================

(This entire section is unchanged from v5.2.)

\subsection{Body-Frame Shape Function}

\begin{equation}
  T(x_b,y_b,z_b) =
  \begin{cases}
    1, & (x_b,y_b,z_b)\in\mathcal{V},\\
    0, & \text{otherwise}.
  \end{cases}
  \label{eq:body_shape}
\end{equation}

\subsection{Body-Frame to World-Frame Transformation}

\begin{align}
  x_w &= x_T(t) + x_b\cos\theta - y_b\sin\theta,
  \label{eq:xw}\\
  y_w &= y_T(t) + x_b\sin\theta + y_b\cos\theta,
  \label{eq:yw}\\
  z_w &= z_T + z_b.
  \label{eq:zw}
\end{align}

\subsection{Shadow Silhouette in the Aperture Plane}
\label{sec:shadow_silhouette}

The effective propagation distance from a body element at
$(x_b,y_b)$ to the receiver plane:
\begin{equation}
  R_{\rm eff}(x_b,y_b) = R_0 - y_w
  = R(t) - x_b\sin\theta - y_b\cos\theta.
  \label{eq:R_eff_general}
\end{equation}

The shadow profile function in aperture-plane coordinates:
\begin{equation}
  F(x',z') =
  \begin{cases}
    1, & \exists\;y_b:\;T(x'/\!\cos\theta + y_b\sin\theta/\!\cos\theta,
    \;y_b,\;z')=1,\\
    0, & \text{otherwise}.
  \end{cases}
  \label{eq:shadow_profile}
\end{equation}

\subsection{Pixel-Dependent Effective Propagation Distance}
\label{sec:pixel_R}

For a thin target ($y_b\approx 0$):
\begin{equation}
  \boxed{R_{\rm eff}(x',t) = R(t) - x'\tan\theta.}
  \label{eq:R_eff_pixel}
\end{equation}

%======================================================================
\section{FSSR Model}
\label{sec:fssr}
%======================================================================

(Unchanged from v5.2.  The FSSR is a property of the target--receiver
geometry, not the array structure within a node.)

\subsection{Continuous FSSR with Pixel-Dependent Propagation}

\begin{equation}
  \boxed{
  \varepsilon_k(t) = \biggl|1
  + \iint F(x',z')\;\frac{i}{R_{\rm eff}(x',t)}\;
  \exp\!\biggl[\frac{i\pi}{R_{\rm eff}(x',t)}\Bigl(
  \bigl(\Delta x_k(t)+x'\bigr)^2
  + \bigl(z_T+z'\bigr)^2\Bigr)\biggr]
  \frac{dx'\,dz'}{\cos\theta}\biggr|^2,}
  \label{eq:fssr_continuous}
\end{equation}
where $\Delta x_k(t)=x_T(t)-\bar{X}_k$.

\subsection{Pixel Discretisation}

$N_{\rm pix}=M_z\times M_x$ square pixels of side~$2\delta$.
Pixel~$(p,q)$ has centre $(X'_q,Z'_p)$ and value
$P_{p,q}\in\{0,1\}$.  Effective propagation distance for column~$q$:
\begin{equation}
  R_{\mathrm{eff},q}(t) = R(t) - X'_q\tan\theta.
  \label{eq:R_eff_q}
\end{equation}

\begin{equation}
  \varepsilon_k(t) \approx \biggl|1
  + \frac{1}{\cos\theta}\sum_{p,q}
  \frac{i}{R_{\mathrm{eff},q}(t)}\;
  F_q^{(k)}(t)\;G_{p,q}(t)\;P_{p,q}\biggr|^2.
  \label{eq:fssr_pixel}
\end{equation}

\subsubsection{Fresnel Coefficients}

$x$-direction:
\begin{equation}
  F_q^{(k)}(t)
  = \sqrt{\frac{R_{\mathrm{eff},q}(t)}{4i}}\,\biggl[
    \erfi\!\biggl(\sqrt{\frac{i\pi}{R_{\mathrm{eff},q}(t)}}
    \bigl(X'_q{+}\delta{+}\Delta x_k(t)\bigr)\biggr)
  - \erfi\!\biggl(\sqrt{\frac{i\pi}{R_{\mathrm{eff},q}(t)}}
    \bigl(X'_q{-}\delta{+}\Delta x_k(t)\bigr)\biggr)\biggr].
  \label{eq:Fq}
\end{equation}

$z$-direction (column-dependent):
\begin{equation}
  G_{p,q}(t)
  = \sqrt{\frac{R_{\mathrm{eff},q}(t)}{4i}}\,\biggl[
    \erfi\!\biggl(\sqrt{\frac{i\pi}{R_{\mathrm{eff},q}(t)}}
    \bigl(Z'_p{+}\delta{+}z_T\bigr)\biggr)
  - \erfi\!\biggl(\sqrt{\frac{i\pi}{R_{\mathrm{eff},q}(t)}}
    \bigl(Z'_p{-}\delta{+}z_T\bigr)\biggr)\biggr].
  \label{eq:Gpq}
\end{equation}

\subsection{Combined Fresnel Coefficient Matrix}

\begin{equation}
  H_{k,p,q}(t) = \frac{i}{\cos\theta\;R_{\mathrm{eff},q}(t)}\;
  F_q^{(k)}(t)\;G_{p,q}(t),
  \label{eq:H_matrix}
\end{equation}
\begin{align}
  S_k(t) &= \sum_{p,q} H_{k,p,q}(t)\,P_{p,q},
  \label{eq:Sk}\\
  \varepsilon_k(t) &= |1 + S_k(t)|^2. \label{eq:eps_from_S}
\end{align}

%======================================================================
\section{Unified Noise Model}
\label{sec:noise}
%======================================================================

\subsection{Complex Signal Noise}
At element $(n_x,n_z)$ of node~$k$:
\begin{equation}
  \eta_{k,n_x,n_z}(t) \sim \mathcal{CN}(0,\,\sigma^2),
  \qquad
  \sigma^2 = 10^{-\mathrm{SNR}_{\mathrm{dB}}/10},
  \label{eq:noise_def}
\end{equation}
referenced to $|A_D|^2=1$.  The noise is i.i.d.\ across all $N^2$
elements.

\subsection{FSSR Measurement via Total Field}
\label{sec:fssr_measurement}

The direct-path amplitude $A_D=1$ is perfectly known.  After
\textbf{2D beamforming} toward broadside at node~$k$ (steering both
$x$ and $z$ dimensions to $(u_x,u_z)=(0,0)$):
\begin{equation}
  \tilde{y}_k(t) = \frac{1}{N^2}\sum_{n_x=0}^{N-1}
  \sum_{n_z=0}^{N-1} x_{k,n_x,n_z}(t)
  \approx 1 + S_k(t) + \tilde{\eta}_k(t),
  \label{eq:beamformed}
\end{equation}
where the beamformed noise is
\begin{equation}
  \tilde{\eta}_k(t) \sim \mathcal{CN}(0,\,\sigma^2/N^2).
  \label{eq:beamformed_noise}
\end{equation}

\paragraph{SNR improvement.}
Compared with v5.2 (where $\tilde{\eta}\sim\mathcal{CN}(0,\sigma^2/N)$),
the 2D array reduces the beamformed noise power by a factor of~$N$,
providing an additional $10\log_{10}N$\,dB of processing gain.

The FSSR is measured as
\begin{equation}
  \boxed{b_k(t) = |\tilde{y}_k(t)|^2
  = |1 + S_k(t) + \tilde{\eta}_k(t)|^2.}
  \label{eq:fssr_measured}
\end{equation}
Expanding:
\begin{equation}
  b_k(t) = \underbrace{|1{+}S_k|^2}_{\varepsilon_k}
  + \underbrace{2\,\mathrm{Re}\!\bigl[(1{+}S_k)\,
  \tilde{\eta}_k^*\bigr]}_{\text{signal-dependent}}
  + \underbrace{|\tilde{\eta}_k|^2}_{\text{quadratic}}.
  \label{eq:fssr_expanded}
\end{equation}

%======================================================================
\section{Zone Classification}
\label{sec:zones}
%======================================================================

(Unchanged from v5.2.)

\subsection{Per-Node Fresnel Parameter}
\begin{equation}
  \xi_k(t) = \frac{|\Delta x_k(t)|}{\sqrt{R(t)}}.
  \label{eq:xi_k}
\end{equation}
\begin{equation}
  \text{Zone~B: } \xi_k(t)>\xi_{\mathrm{th}},
  \qquad
  \text{Zone~A: } \xi_k(t)\le\xi_{\mathrm{th}}.
  \label{eq:zone_def}
\end{equation}

\subsection{Sequential Zone Entry}
Node~$k$ enters Zone~A around time
$t_k^*\approx(\bar{X}_k-x_0)/v_x$.

%======================================================================
\section{Zone~B: Parameter Estimation with 2D Arrays}
\label{sec:zone_b}
%======================================================================

When $\xi_k\gg 1$, the target appears as a point source.  The 2D
array enables estimation of \emph{both} the azimuth direction
cosine~$\ell_x$ and the elevation direction cosine~$\ell_z$,
allowing direct estimation of $z_T$ from Zone~B data.

The unknowns are $v_x,v_y,x_0,R_0,z_T$
(or equivalently $v,\theta,x_0,R_0,z_T$).

\subsection{Received Signal Model}
\label{sec:zone_b_signal}

At element $(n_x,n_z)$ of node~$k$:
\begin{equation}
  \boxed{
  x_{k,n_x,n_z}(t) = A_D + G_T(t)\,e^{j\psi_k(t)}\,
  e^{j\pi n_x\,\ell_x^{(k)}(t)}\,
  e^{j\pi n_z\,\ell_z^{(k)}(t)}
  + \eta_{k,n_x,n_z}(t),}
  \label{eq:rx_signal}
\end{equation}
where:
\begin{itemize}
  \item $A_D=1$;
  \item $G_T(t) = 2\sqrt{\pi}\,A_\Sigma / R(t)$: scattered
    amplitude;
  \item $\psi_k(t) = \pi[\Delta x_k(t)^2{+}z_T^2]/R(t)$:
    propagation phase;
  \item $\ell_x^{(k)}(t) = \Delta x_k(t)/D_k(t)$: $x$-direction
    cosine;
  \item $\ell_z^{(k)}(t) = z_T/D_k(t)$: $z$-direction cosine
    (\textbf{new});
  \item $D_k(t) = \sqrt{\Delta x_k(t)^2+R(t)^2+z_T^2}$: range;
  \item $\eta_{k,n_x,n_z}\sim\mathcal{CN}(0,\sigma^2)$.
\end{itemize}

\paragraph{Key difference from v5.2.}
The phase progression across the $z$-dimension of the array,
$\exp(j\pi n_z\,\ell_z^{(k)})$, is the new term that enables
elevation DOA estimation.  This phase is proportional to $z_T/D_k$
and varies slowly with time (since $z_T$ is constant and $D_k$
changes slowly in Zone~B).

\subsection{2D Steering Vector}
\label{sec:2d_steering}

Define the 2D steering vector for direction cosines $(\ell_x,\ell_z)$:
\begin{equation}
  \mathbf{a}(\ell_x,\ell_z) = \mathbf{a}_x(\ell_x)\otimes
  \mathbf{a}_z(\ell_z) \;\in\;\C^{N^2},
  \label{eq:steering_2d}
\end{equation}
where $\otimes$ denotes the Kronecker product, and
\begin{align}
  \mathbf{a}_x(u) &= [1,\,e^{j\pi u},\,\ldots,\,
  e^{j\pi(N-1)u}]^T \;\in\;\C^N,
  \label{eq:ax}\\
  \mathbf{a}_z(w) &= [1,\,e^{j\pi w},\,\ldots,\,
  e^{j\pi(N-1)w}]^T \;\in\;\C^N.
  \label{eq:az}
\end{align}

The $N^2\times 1$ received signal vector at node~$k$, formed by
stacking columns (i.e.\ vectorising over $(n_x,n_z)$):
\begin{equation}
  \mathbf{x}_k(t) = A_D\,\mathbf{1}_{N^2}
  + G_T(t)\,e^{j\psi_k(t)}\,
  \mathbf{a}\bigl(\ell_x^{(k)}(t),\,\ell_z^{(k)}(t)\bigr)
  + \bm{\eta}_k(t).
  \label{eq:rx_vector}
\end{equation}

\subsection{The Time-Varying DOA Problem}
\label{sec:tv_doa}

The azimuth DOA $\ell_x^{(k)}(t)$ varies significantly with time
(same as v5.2).  The elevation DOA $\ell_z^{(k)}(t)=z_T/D_k(t)$
varies much more slowly because $z_T$ is constant and $D_k$ changes
slowly in Zone~B ($D_k\gg\Delta x_k$).

\subsubsection{Azimuth DOA Variation (Review from v5.2)}

The standard spatial covariance
$\hat{\mathbf{R}}_k=(1/M)\sum_m\mathbf{x}_k(t_m)\mathbf{x}_k^H(t_m)$
suffers from DOA smearing in the $x$-dimension due to the
time-varying $\ell_x$.  The maximum number of uncompensated snapshots:
\begin{equation}
  M_{\max}^{\mathrm{(unc)}} \approx
  \frac{2}{(N{-}1)\,|\beta_x|\,\Delta t},
  \label{eq:Mmax_uncomp}
\end{equation}
where $\beta_x$ is the azimuth DOA rate.

\subsubsection{Elevation DOA Variation}

The elevation DOA rate is
\begin{equation}
  \dot{\ell}_z^{(k)}(t) = -\frac{z_T}{D_k^2}\,\dot{D}_k(t),
  \label{eq:lz_rate}
\end{equation}
where $\dot{D}_k = (v_x\Delta x_k + v_y R)/D_k$ is the range rate.
Since $z_T/D_k\ll 1$ and $\dot{D}_k/D_k$ is small in Zone~B, the
elevation DOA variation is typically much smaller than the azimuth
variation:
\begin{equation}
  |\dot{\ell}_z| \ll |\dot{\ell}_x|.
  \label{eq:lz_slow}
\end{equation}
This means the $z$-dimension of the covariance matrix is minimally
affected by DOA smearing, even for moderately large~$M$.

\subsubsection{Phase-Compensated 2D Spatial Covariance}
\label{sec:phase_comp}

Given an estimate $\hat{\beta}_x$ of the azimuth DOA rate, define
compensated snapshots by dechirping only the $x$-dimension phase:
\begin{equation}
  \boxed{
  \tilde{x}_{k,n_x,n_z}(t_m) = x_{k,n_x,n_z}(t_m)\;
  e^{-j\pi\,n_x\;\!\hat{\beta}_x\,(t_m - t_{\mathrm{ref}})}.}
  \label{eq:phase_comp}
\end{equation}

Note: no compensation is needed in the $z$-dimension since
$\ell_z$ varies negligibly.

The compensated $N^2\times N^2$ spatial covariance:
\begin{equation}
  \tilde{\mathbf{R}}_k
  = \frac{1}{M}\sum_{m=1}^{M}
  \tilde{\mathbf{x}}_k(t_m)\,\tilde{\mathbf{x}}_k^H(t_m).
  \label{eq:cov_compensated}
\end{equation}

If $\hat{\beta}_x=\beta_x$, the azimuth DOA is frozen at
$\ell_x(t_{\mathrm{ref}})$ for all snapshots.

\subsubsection{Optimal~$M$ (unchanged from v5.2 except array size)}
\label{sec:optimal_M}

The azimuth-dimension constraint on~$M$ remains:

\paragraph{(a) Rate estimation error $\Delta\beta_x$:}
$|\Delta\beta_x|\,M\,\Delta t \le c_1/N$
$\;\Longrightarrow\;M \le c_1/(N|\Delta\beta_x|\Delta t)$.

\paragraph{(b) Quadratic azimuth DOA variation~$\gamma_x$:}
\begin{equation}
  \gamma_x^{(k)} \approx \frac{v_x\,v_y}{R_0^2}
  + \frac{v_y^2\,(x_0{-}\bar{X}_k)}{R_0^3}.
  \label{eq:gamma_x}
\end{equation}
$|\gamma_x|(M\Delta t)^2/4 \le c_2/N$
$\;\Longrightarrow\;M \le 2\sqrt{c_2/N}\,/\,(|\gamma_x|^{1/2}\Delta t)$.

\paragraph{Combined optimal~$M$:}
\begin{equation}
  \boxed{
  M_{\mathrm{opt}} = \min\!\left(
    \frac{c_1}{N\,|\Delta\beta_x|\,\Delta t},\;
    \frac{2}{\sqrt{N\,|\gamma_x|}\;\!\Delta t},\;
    M_{\mathrm{avail}}\right).}
  \label{eq:M_opt}
\end{equation}

\paragraph{Iterative refinement:}
\begin{enumerate}
  \item \textbf{Pass~1} (no compensation): $M\le M_{\max}^{(\mathrm{unc})}$;
    rough~$\hat{\beta}_x$.
  \item \textbf{Pass~2+}: compensate with~$\hat{\beta}_x$, enlarge $M$
    to~$M_{\mathrm{opt}}$, re-estimate DOA.
\end{enumerate}

\subsection{2D DOA Estimation}
\label{sec:2d_doa}

The Kronecker structure of the 2D steering
vector~\eqref{eq:steering_2d} allows efficient separable or joint
DOA estimation.

\subsubsection{Method 1: Joint 2D MUSIC}
\label{sec:2d_music}

Apply eigendecomposition to
$\tilde{\mathbf{R}}_k\in\C^{N^2\times N^2}$; partition into signal
and noise subspaces $\mathbf{E}_s$, $\mathbf{E}_n$.  The 2D MUSIC
pseudo-spectrum:
\begin{equation}
  P_{\mathrm{MUSIC}}^{(2D)}(\ell_x,\ell_z)
  = \frac{1}{\mathbf{a}^H(\ell_x,\ell_z)\,
  \mathbf{E}_n\mathbf{E}_n^H\,\mathbf{a}(\ell_x,\ell_z)}.
  \label{eq:music_2d}
\end{equation}
The peak location gives $(\hat{\ell}_x,\hat{\ell}_z)$.

\paragraph{Computational cost.}
The eigendecomposition of $\tilde{\mathbf{R}}_k\in\C^{N^2\times N^2}$
has cost $O(N^6)$.  For moderate $N$ ($\le 16$), this is tractable.
For larger $N$, use the separable methods below.

\subsubsection{Method 2: Separable Estimation via Row and Column
Covariances}
\label{sec:separable_doa}

Exploit the Kronecker structure
$\mathbf{R}_{\mathrm{sig}}=\sigma_s^2\,
\mathbf{a}_x\mathbf{a}_x^H\otimes\mathbf{a}_z\mathbf{a}_z^H$.

\paragraph{Step 1: Azimuth DOA.}
Form the \textbf{row-averaged covariance} by summing the $N\times N$
diagonal blocks of $\tilde{\mathbf{R}}_k$:
\begin{equation}
  \mathbf{R}_x = \frac{1}{N}\sum_{m=0}^{N-1}
  [\tilde{\mathbf{R}}_k]_{mN:(m+1)N,\;mN:(m+1)N}
  \;\in\;\C^{N\times N}.
  \label{eq:Rx}
\end{equation}
Apply MUSIC or LS-ESPRIT to $\mathbf{R}_x$ to estimate $\hat{\ell}_x$.

\paragraph{Step 2: Elevation DOA.}
Form the \textbf{column-averaged covariance} by extracting the
$N\times N$ blocks at positions $(m_1,m_2)$ along the
$x$-dimension and summing:
\begin{equation}
  [\mathbf{R}_z]_{n_{z_1},n_{z_2}}
  = \frac{1}{N}\sum_{m=0}^{N-1}
  [\tilde{\mathbf{R}}_k]_{mN+n_{z_1},\;mN+n_{z_2}}
  \;\in\;\C^{N\times N}.
  \label{eq:Rz}
\end{equation}
Apply MUSIC or LS-ESPRIT to $\mathbf{R}_z$ to estimate $\hat{\ell}_z$.

\paragraph{Advantage.}
Each step requires only an $N\times N$ eigendecomposition ($O(N^3)$),
reducing cost from $O(N^6)$ to $O(N^3)$.  The averaging also improves
the effective number of snapshots for each dimension.

\subsubsection{Method 3: 2D LS-ESPRIT}
\label{sec:2d_esprit}

Define shift-invariance structures for both dimensions:

\paragraph{Azimuth ($x$-dimension):}
Subarrays formed by selecting rows $0,\ldots,N{-}2$ vs.\
$1,\ldots,N{-}1$ of the $x$-dimension.  For the Kronecker-structured
signal subspace $\mathbf{E}_s$:
\begin{equation}
  \mathbf{E}_{x,1} = \mathbf{J}_{x,1}\,\mathbf{E}_s, \qquad
  \mathbf{E}_{x,2} = \mathbf{J}_{x,2}\,\mathbf{E}_s,
  \label{eq:esprit_x}
\end{equation}
where $\mathbf{J}_{x,1}=\mathbf{I}_{N-1,N}^{(\mathrm{top})}
\otimes\mathbf{I}_N$ and
$\mathbf{J}_{x,2}=\mathbf{I}_{N-1,N}^{(\mathrm{bot})}
\otimes\mathbf{I}_N$ are the selection matrices.  Then
$\hat{\ell}_x=\angle\lambda_x/\pi$ from the eigenvalue of
$\mathbf{E}_{x,1}^\dagger\mathbf{E}_{x,2}$.

\paragraph{Elevation ($z$-dimension):}
Subarrays by selecting columns $0,\ldots,N{-}2$ vs.\
$1,\ldots,N{-}1$ of the $z$-dimension:
\begin{equation}
  \mathbf{E}_{z,1} = \mathbf{J}_{z,1}\,\mathbf{E}_s, \qquad
  \mathbf{E}_{z,2} = \mathbf{J}_{z,2}\,\mathbf{E}_s,
  \label{eq:esprit_z}
\end{equation}
where $\mathbf{J}_{z,1}=\mathbf{I}_N\otimes
\mathbf{I}_{N-1,N}^{(\mathrm{top})}$ and
$\mathbf{J}_{z,2}=\mathbf{I}_N\otimes
\mathbf{I}_{N-1,N}^{(\mathrm{bot})}$.  Then
$\hat{\ell}_z=\angle\lambda_z/\pi$.

\subsection{Altitude Estimation from Elevation DOA}
\label{sec:zT_from_doa}

Given $\hat{\ell}_z^{(k)}(t)$ from the 2D array at node~$k$:

\subsubsection{Far-Field Approximation}

For $D_k\approx R(t)$:
\begin{equation}
  \ell_z^{(k)}(t) \approx \frac{z_T}{R(t)}.
  \label{eq:lz_farfield}
\end{equation}
This is \textbf{independent of $k$} in the far-field limit (since
$z_T$ and $R(t)$ are the same for all nodes).

Estimate:
\begin{equation}
  \boxed{\hat{z}_T = \hat{\ell}_z\;\hat{R}(t),}
  \label{eq:zT_from_lz}
\end{equation}
where $\hat{R}(t)$ is estimated from the azimuth DOA and Doppler
analysis (Section~\ref{sec:joint_zone_b}).

\subsubsection{Exact Inversion}

Using $D_k=\sqrt{\Delta x_k^2+R(t)^2+z_T^2}$:
\begin{equation}
  z_T = \frac{\hat{\ell}_z\sqrt{\Delta x_k^2+R^2}}
  {\sqrt{1-\hat{\ell}_z^2}}.
  \label{eq:zT_exact}
\end{equation}

\subsubsection{Time-Averaged $z_T$ Estimate}

Since $z_T$ is constant, average over $L$ time steps and $K$ nodes:
\begin{equation}
  \hat{z}_T^{(\mathrm{B})} = \frac{1}{KL}\sum_{k=0}^{K-1}
  \sum_{\ell=1}^{L}
  \frac{\hat{\ell}_z^{(k)}(t_\ell)\,\hat{R}(t_\ell)}
  {\sqrt{1-[\hat{\ell}_z^{(k)}(t_\ell)]^2}}\;\!
  \sqrt{1+\bigl[\hat{\ell}_x^{(k)}(t_\ell)\bigr]^2
  /\bigl(1-[\hat{\ell}_x^{(k)}(t_\ell)]^2\bigr)}.
  \label{eq:zT_averaged}
\end{equation}
In the far-field limit, this simplifies to
$\hat{z}_T^{(\mathrm{B})}\approx(1/KL)\sum_{k,\ell}
\hat{\ell}_z^{(k)}(t_\ell)\,\hat{R}(t_\ell)$.

\subsubsection{Achievable Precision}

The Cram\'er--Rao bound for the elevation direction cosine from an
$N$-element ULA with SNR per element~$\rho$ and $M$ snapshots is:
\begin{equation}
  \mathrm{var}(\hat{\ell}_z)
  \ge \frac{6}{\pi^2\,\rho\,M\,N(N^2{-}1)}.
  \label{eq:crb_lz}
\end{equation}
The corresponding altitude precision (far-field):
\begin{equation}
  \mathrm{std}(\hat{z}_T) \approx R
  \sqrt{\frac{6}{\pi^2\,\rho\,M\,N(N^2{-}1)}}.
  \label{eq:std_zT}
\end{equation}
For $N=16$, $\rho=100$ (20\,dB), $M=100$, $R=10^4$:
$\mathrm{std}(\hat{z}_T)\approx 1.7\lambda$, sufficient for a
coarse altitude estimate that can be refined in Zone~A.

\subsection{Azimuth DOA Time Series and Nonlinear Model}
\label{sec:doa_nonlinear}

(Unchanged from v5.2.)
In the far field:
\begin{equation}
  \ell_x^{(k)}(t) \approx \frac{x_0+v_x t-\bar{X}_k}{R_0-v_y t}.
  \label{eq:doa_ratio}
\end{equation}
Taylor expansion:
\begin{align}
  \ell_x^{(k)}(t) &\approx \alpha_k + \beta_k\,t + \gamma_k\,t^2
  + \cdots, \label{eq:doa_taylor}
\end{align}
\begin{align}
  \alpha_k &= \frac{x_0{-}\bar{X}_k}{R_0},
  \label{eq:alpha_k}\\
  \beta_k &= \frac{v_x}{R_0}
  + \frac{v_y(x_0{-}\bar{X}_k)}{R_0^2},
  \label{eq:beta_k}\\
  \gamma_k &= \frac{v_x v_y}{R_0^2}
  + \frac{v_y^2(x_0{-}\bar{X}_k)}{R_0^3}.
  \label{eq:gamma_k}
\end{align}

Node-dependent DOA rate:
\begin{equation}
  \beta_k = \underbrace{\frac{v_x}{R_0}
  + \frac{v_y\,x_0}{R_0^2}}_a
  - \underbrace{\frac{v_y}{R_0^2}}_b\;\bar{X}_k
  \;\triangleq\; a - b\,\bar{X}_k.
  \label{eq:beta_linear_in_Xk}
\end{equation}

\subsection{Elevation DOA Time Series}
\label{sec:elev_doa_timeseries}

The elevation direction cosine in the far field:
\begin{equation}
  \ell_z^{(k)}(t) \approx \frac{z_T}{R(t)}
  = \frac{z_T}{R_0 - v_y\,t}.
  \label{eq:lz_time}
\end{equation}
For $|v_y t|\ll R_0$:
\begin{equation}
  \ell_z^{(k)}(t) \approx \frac{z_T}{R_0}
  + \frac{z_T\,v_y}{R_0^2}\,t + \cdots
  \;\triangleq\; \alpha_z + \beta_z\,t,
  \label{eq:lz_taylor}
\end{equation}
with
\begin{equation}
  \alpha_z = \frac{z_T}{R_0},\qquad
  \beta_z = \frac{z_T\,v_y}{R_0^2}.
  \label{eq:lz_coeffs}
\end{equation}

\paragraph{Key observations:}
\begin{enumerate}
  \item $\alpha_z$ directly gives $z_T/R_0$; combined with
    $\hat{R}_0$ from azimuth analysis, this yields $\hat{z}_T$.
  \item $\beta_z$ is small when $v_y$ is small (near-perpendicular
    crossing) or when $z_T$ is small.  In either case, $\ell_z$
    varies slowly, confirming that $z$-dimension phase compensation
    is unnecessary.
  \item For $\theta=0$: $\beta_z=0$ and $\ell_z$ is strictly
    constant.
\end{enumerate}

\subsection{Doppler Analysis}
\label{sec:doppler}

(Unchanged from v5.2.)
The propagation phase:
$\psi_k(t)=\pi[\Delta x_k(t)^2+z_T^2]/R(t)$.  Instantaneous
Doppler:
\begin{equation}
  f_D^{(k)}(t) = \frac{1}{2\pi}\frac{d\psi_k}{dt}
  = \frac{v_x\,\Delta x_k(t)}{R(t)}
  + \frac{v_y[\Delta x_k(t)^2+z_T^2]}{2\,R(t)^2}.
  \label{eq:doppler_full}
\end{equation}

The Doppler is estimated via STFT on the 2D-steered beam.  With
estimated direction cosines $(\hat{\ell}_x,\hat{\ell}_z)$:
\begin{equation}
  y_k(t) = \frac{1}{N^2}\sum_{n_x,n_z}
  x_{k,n_x,n_z}(t)\;e^{-j\pi(n_x\hat{\ell}_x+n_z\hat{\ell}_z)}.
  \label{eq:steered_beam}
\end{equation}

\paragraph{Note.}
In v5.2, the steered beam was
$(1/N)\sum_{n_x}x_{k,n_x}e^{-j\pi n_x\hat{\ell}_x}$ (1D steering).
The 2D steering~\eqref{eq:steered_beam} provides an additional
$10\log_{10}N$\,dB of gain by also coherently combining the
$z$-dimension elements toward the target.

\subsection{Joint Parameter Estimation}
\label{sec:joint_zone_b}

\subsubsection{Step 1: $v_y/R_0^2$ from azimuth DOA-rate variation}
(Unchanged.)
From~\eqref{eq:beta_linear_in_Xk}:
$\hat{b}=-v_y/R_0^2$, $\hat{a}=v_x/R_0+v_y x_0/R_0^2$.

\subsubsection{Step 2: Combining with Doppler}
(Unchanged.)
For $\theta\approx 0$:
$|v_x|=|\hat{\dot{f}}|/|\hat{\beta}|$,\;
$R_0=|v_x|/|\hat{\beta}|$.

For $\theta\ne 0$: nonlinear least-squares fit.  Initialisation:
\begin{equation}
  \hat{R}_0 \approx
  \frac{|\hat{\dot{f}}_{\mathrm{avg}}|}
  {\hat{\beta}_{\mathrm{avg}}^2},\qquad
  \hat{v}_y = -\hat{b}\,\hat{R}_0^2,\qquad
  \hat{v}_x = (\hat{a}-\hat{b}\,\hat{x}_0)\hat{R}_0,\qquad
  \hat{x}_0 = \hat{\alpha}_{\mathrm{avg}}\hat{R}_0
  +\bar{X}_{\mathrm{avg}}.
  \label{eq:joint_init}
\end{equation}

\subsubsection{Step 3: $z_T$ from elevation DOA (\textbf{new in v6})}

Using the elevation DOA estimates from
Section~\ref{sec:zT_from_doa}:
\begin{equation}
  \boxed{
  \hat{z}_T^{(\mathrm{B})} = \hat{\alpha}_z\;\hat{R}_0
  = \frac{1}{KL}\sum_{k,\ell}\hat{\ell}_z^{(k)}(t_\ell)\;\hat{R}_0.}
  \label{eq:zT_zone_b}
\end{equation}

\paragraph{Cross-check via elevation DOA rate.}
If $v_y\ne 0$, the elevation DOA rate provides a consistency check:
\begin{equation}
  \hat{z}_T = \frac{\hat{\beta}_z\,\hat{R}_0^2}{\hat{v}_y}.
  \label{eq:zT_crosscheck}
\end{equation}

%======================================================================
\section{Zone~A: Gradient Descent with TV Denoising}
\label{sec:zone_a}
%======================================================================

(The Zone~A formulation is unchanged from v5.2.  The only difference
is that the initial $z_T$ estimate is now $\hat{z}_T^{(\mathrm{B})}$
from the elevation DOA, rather than $z_T=0$.)

\subsection{Problem Formulation}

\begin{equation}
  \hat{\mathbf{P}} = \argmin_{\mathbf{P}}
  \bigl\{\mathcal{F}(\mathbf{P}) + \tau\,\mathcal{R}(\mathbf{P})\bigr\},
  \label{eq:zone_a_problem}
\end{equation}
\begin{equation}
  \mathcal{F}(\mathbf{P})
  = \sum_{k,\ell}
  \bigl[b_{k,\ell} - \varepsilon_k(t_\ell;\mathbf{P})\bigr]^2,
  \label{eq:fidelity}
\end{equation}
\begin{equation}
  \mathcal{R}(\mathbf{P}) = \TV(\mathbf{P})
  = \sum_{p,q}\sqrt{(P_{p,q}{-}P_{p+1,q})^2
  +(P_{p,q}{-}P_{p,q+1})^2}.
  \label{eq:tv_reg}
\end{equation}

\subsection{Gradient of the Fidelity Function}
\label{sec:gradient}

Using the combined coefficient~\eqref{eq:H_matrix}:
\begin{equation}
  \boxed{
  \frac{\partial\mathcal{F}_{k,\ell}}{\partial P_{p,q}}
  = 4[\varepsilon_k{-}b_{k,\ell}]\,
  \mathrm{Re}\!\bigl[H_{k,p,q}(t_\ell)\,
  \overline{(1{+}S_{k,\ell})}\bigr].}
  \label{eq:grad_fidelity}
\end{equation}

\subsection{Algorithm 1: Shadow Profile Retrieval}
\label{sec:alg1}

\begin{algorithm}[H]
\caption{Gradient descent with TV denoising for shadow profile
retrieval}
\label{alg:profile_retrieval}
\begin{algorithmic}[1]
\Require FSSR measurements $\{b_{k,\ell}\}$; coefficient matrices
$\{\mathbf{H}_k(t_\ell)\}$; step size $\gamma>0$;
TV parameter $\tau>0$; denoising iterations~$N_D$
\State Initialise: $\hat{\mathbf{P}}_0\leftarrow\mathbf{0}$,
$\mathbf{s}_0\leftarrow\hat{\mathbf{P}}_0$, $q_0\leftarrow 1$
\For{$t=1,\ldots,N_S$}
  \State $\gamma_t \leftarrow \gamma/\sqrt{t}$
  \State $\mathbf{z}_t \leftarrow \mathbf{s}_{t-1}
    - \frac{\gamma_t}{L_{\mathrm{tot}}}
    \sum_{k,\ell}\nabla_{\mathbf{P}}\mathcal{F}_{k,\ell}
    (\mathbf{s}_{t-1})$
  \Comment{$L_{\mathrm{tot}}=\sum_k L_k$}
  \State $\hat{\mathbf{P}}_t \leftarrow
    \mathrm{prox}_{\mathcal{R}}(\mathbf{z}_t,\gamma_t\tau)$
  \Comment{Algorithm~\ref{alg:tv_denoise}}
  \State $q_t \leftarrow (1{+}\sqrt{1{+}4q_{t-1}^2})/2$
  \State $\mathbf{s}_t \leftarrow \hat{\mathbf{P}}_t
    + \frac{q_{t-1}-1}{q_t}
    (\hat{\mathbf{P}}_t{-}\hat{\mathbf{P}}_{t-1})$
\EndFor
\State \Return $\hat{\mathbf{P}}_{N_S}$
\end{algorithmic}
\end{algorithm}

\subsection{Algorithm 2: TV Denoising Proximal Operator}
\label{sec:alg2}

\begin{equation}
  \mathrm{prox}_{\mathcal{R}}(\mathbf{z},\delta)
  = \argmin_{\mathbf{P}}\Bigl\{
  \tfrac{1}{2\delta}\|\mathbf{P}{-}\mathbf{z}\|^2
  + \TV(\mathbf{P})\Bigr\}
  \quad\text{s.t.}\quad 0\le P_{p,q}\le 1.
  \label{eq:prox_problem}
\end{equation}

\paragraph{Operators.}
\begin{itemize}
  \item \textbf{Discrete gradient}
    $\mathcal{L}^T$:
    $p_{i,j}=m_{i,j}{-}m_{i+1,j}$,\;
    $q_{i,j}=m_{i,j}{-}m_{i,j+1}$.
  \item \textbf{Discrete divergence}
    $\mathcal{L}$:
    $[\mathcal{L}(\mathbf{p},\mathbf{q})]_{i,j}
    = p_{i,j}{+}q_{i,j}{-}p_{i-1,j}{-}q_{i,j-1}$.
  \item \textbf{Clipping}
    $\mathcal{H}_C$:
    $[\mathcal{H}_C(\mathbf{m})]_{i,j}=\min(1,\max(0,m_{i,j}))$.
  \item \textbf{Ball projection}
    $\mathcal{H}_B$:
    $r_{i,j}=p_{i,j}/\max(1,\sqrt{p_{i,j}^2{+}q_{i,j}^2})$,\;
    $s_{i,j}=q_{i,j}/\max(1,\sqrt{p_{i,j}^2{+}q_{i,j}^2})$.
\end{itemize}

\begin{algorithm}[H]
\caption{TV denoising:
$\mathrm{prox}_{\mathcal{R}}(\mathbf{z},\delta)$}
\label{alg:tv_denoise}
\begin{algorithmic}[1]
\Require Pixel image $\mathbf{z}\in\R^{M_z\times M_x}$,
regularisation $\delta$
\State $(\mathbf{p}_0,\mathbf{q}_0)\leftarrow(\mathbf{0},\mathbf{0})$;
$(\mathbf{r}_1,\mathbf{s}_1)\leftarrow(\mathbf{0},\mathbf{0})$;
$w_1\leftarrow 1$
\For{$t=1,\ldots,N_D$}
  \State $(\mathbf{p}_t,\mathbf{q}_t) \leftarrow
    \mathcal{H}_B\!\bigl[(\mathbf{r}_t,\mathbf{s}_t)
    + \tfrac{1}{8\delta}\,\mathcal{L}^T\!\bigl(
    \mathcal{H}_C[\mathbf{z}{-}\delta\,
    \mathcal{L}(\mathbf{r}_t,\mathbf{s}_t)]\bigr)\bigr]$
  \State $w_{t+1}\leftarrow(1{+}\sqrt{1{+}4w_t^2})/2$
  \State $(\mathbf{r}_{t+1},\mathbf{s}_{t+1})\leftarrow
    (\mathbf{p}_t,\mathbf{q}_t)
    +\tfrac{w_t-1}{w_{t+1}}
    [(\mathbf{p}_t{-}\mathbf{p}_{t-1},\,
    \mathbf{q}_t{-}\mathbf{q}_{t-1})]$
\EndFor
\State $\hat{\mathbf{x}}\leftarrow
  \mathcal{H}_C[\mathbf{z}{-}\delta\,
  \mathcal{L}(\mathbf{r}_{N_D},\mathbf{s}_{N_D})]$
\State \Return $\hat{\mathbf{x}}$
\end{algorithmic}
\end{algorithm}

\subsection{Algorithm 3: Joint Motion Parameter and Profile Estimation}
\label{sec:alg3}

The outer loop optimises $\bm{\theta}=(v_x,v_y,z_T,R_0)$ around the
profile retrieval, \textbf{initialised with the Zone~B estimates
including $\hat{z}_T^{(\mathrm{B})}$} from the elevation DOA.

The cost function:
\begin{equation}
  C(\hat{\mathbf{P}},\bm{\theta})
  = \sum_{k,\ell}\bigl[b_{k,\ell}
  - \varepsilon_k(t_\ell;\hat{\mathbf{P}},\bm{\theta})\bigr]^2.
  \label{eq:cost_joint}
\end{equation}

\begin{algorithm}[H]
\caption{Joint shadow profile and motion parameter estimation}
\label{alg:joint_estimation}
\begin{algorithmic}[1]
\Require FSSR measurements $\{b_{k,\ell}\}$; initial estimates
$\hat{v}_x^{(0)},\hat{v}_y^{(0)},
\hat{z}_T^{(0)}=\hat{z}_T^{(\mathrm{B})},\hat{R}_0^{(0)},
\hat{x}_0$ (fixed)
\State $q_0\leftarrow 1$;
$\bm{\theta}^{(0)}\leftarrow
(\hat{v}_x^{(0)},\hat{v}_y^{(0)},\hat{z}_T^{(0)},\hat{R}_0^{(0)})$
\For{$n=1,\ldots,N_V$}
  \State Compute trajectories and
    $R_{\mathrm{eff},q}(t_\ell)=R(t_\ell)-X'_q\tan\theta^{(n-1)}$
    \Comment{$\theta=\arctan(v_y/v_x)$}
  \State Build coefficient matrices
    $\mathbf{H}_k(t_\ell)$ using~\eqref{eq:H_matrix}
  \State $\hat{\mathbf{P}}_n \leftarrow
    \text{Algorithm~\ref{alg:profile_retrieval}}
    (\{b_{k,\ell}\},\bm{\theta}^{(n-1)})$
  \State Compute $C$ and numerical gradients
    $\partial C/\partial\theta_i\approx
    [C(\bm{\theta}{+}\Delta\theta_i\mathbf{e}_i)-C(\bm{\theta})]
    /\Delta\theta_i$
  \State $\gamma_i^{(n)}\leftarrow\gamma_i/n$;\quad
    $\tilde{\theta}_i \leftarrow \theta_i^{(n-1)}
    - \gamma_i^{(n)}\,\partial C/\partial\theta_i$
  \State Nesterov:
    $q_n\leftarrow(1{+}\sqrt{1{+}4q_{n-1}^2})/2$;\quad
    $\theta_i^{(n)}\leftarrow\tilde{\theta}_i
    +\tfrac{q_{n-1}-1}{q_n}(\tilde{\theta}_i{-}\theta_i^{(n-1)})$
  \State Enforce $R_0^{(n)}\ge R_{\min}$
\EndFor
\State Final retrieval with converged~$\bm{\theta}$; hard-threshold at
$0.5$
\State \Return
  $\hat{\mathbf{P}},\hat{v}_x,\hat{v}_y,\hat{z}_T,\hat{R}_0$
\end{algorithmic}
\end{algorithm}

%======================================================================
\section{Scale Ambiguity Analysis}
\label{sec:ambiguity}
%======================================================================

(Unchanged from v5.2.)

\subsection{Single-Receiver Ambiguity}
Under the scaling
$(v_x,v_y,z_T,x_0)\to(\mu v_x,\mu v_y,\mu z_T,\mu x_0)$,
$R_0\to\mu^2 R_0$,
$F'(x',z')=F(x'/\mu,z'/\mu)$,
the FSSR is unchanged.

\subsection{Ambiguity Breaking with Multiple Nodes}
With receivers at known, unscaled positions~$\bar{X}_k$:
$\Delta x_k\to\mu x_T-\bar{X}_k\ne\mu\Delta x_k$ for
$\bar{X}_k\ne 0$.

%======================================================================
\section{2D DOA Estimation Algorithms}
\label{sec:doa_algs}
%======================================================================

\subsection{2D MUSIC}
\label{sec:music_2d_alg}

Applied to the phase-compensated covariance~\eqref{eq:cov_compensated}
of size $N^2\times N^2$.  Noise subspace~$\mathbf{E}_n$;
pseudo-spectrum:
\begin{equation}
  P_{\mathrm{MUSIC}}^{(2D)}(\ell_x,\ell_z)
  = \frac{1}{\mathbf{a}^H(\ell_x,\ell_z)\,
  \mathbf{E}_n\mathbf{E}_n^H\,\mathbf{a}(\ell_x,\ell_z)},
  \label{eq:music_2d_alg}
\end{equation}
with $\mathbf{a}(\ell_x,\ell_z)=\mathbf{a}_x(\ell_x)\otimes
\mathbf{a}_z(\ell_z)$.

\subsection{Separable MUSIC}
\label{sec:sep_music}

\paragraph{Azimuth:} MUSIC on
$\mathbf{R}_x\in\C^{N\times N}$~\eqref{eq:Rx}:
\begin{equation}
  P_x(\ell_x) = \frac{1}{\mathbf{a}_x^H(\ell_x)\,
  \mathbf{E}_{n,x}\mathbf{E}_{n,x}^H\,\mathbf{a}_x(\ell_x)}.
  \label{eq:music_x}
\end{equation}

\paragraph{Elevation:} MUSIC on
$\mathbf{R}_z\in\C^{N\times N}$~\eqref{eq:Rz}:
\begin{equation}
  P_z(\ell_z) = \frac{1}{\mathbf{a}_z^H(\ell_z)\,
  \mathbf{E}_{n,z}\mathbf{E}_{n,z}^H\,\mathbf{a}_z(\ell_z)}.
  \label{eq:music_z}
\end{equation}

\subsection{LS-ESPRIT (1D and 2D)}
\label{sec:esprit}

\paragraph{1D LS-ESPRIT (per dimension).}
Signal subspace $\mathbf{E}_s$; subarrays
$\mathbf{E}_1=\mathbf{E}_s(1{:}N{-}1,:)$,
$\mathbf{E}_2=\mathbf{E}_s(2{:}N,:)$;
$\bm{\Phi}=\mathbf{E}_1^\dagger\mathbf{E}_2$;
$\hat{\ell}=\angle\lambda_1/\pi$.

Applied separately to $\mathbf{R}_x$ (for $\hat{\ell}_x$) and
$\mathbf{R}_z$ (for $\hat{\ell}_z$).

\paragraph{2D LS-ESPRIT.}
Applied to the full $N^2\times N^2$ covariance with selection
matrices as described in Section~\ref{sec:2d_esprit}.

\subsection{Exact DOA-to-Position Inversion}
\begin{align}
  \Delta x_k &= \frac{R(t)\,\hat{\ell}_x}
  {\sqrt{1-\hat{\ell}_x^2-\hat{\ell}_z^2}},
  \label{eq:exact_inv_x}\\
  z_T &= \frac{R(t)\,\hat{\ell}_z}
  {\sqrt{1-\hat{\ell}_x^2-\hat{\ell}_z^2}}.
  \label{eq:exact_inv_z}
\end{align}

%======================================================================
\section{Simulation Parameters}
\label{sec:params}
%======================================================================

\begin{table}[H]
\centering
\caption{Default simulation parameters (changes from v5.2
highlighted)}
\label{tab:params}
\begin{tabular}{@{}lcll@{}}
\toprule
\textbf{Symbol} & \textbf{Default} & \textbf{Description}
  & \textbf{Change} \\
\midrule
$K$ & 4 & Receiver nodes & --- \\
$N$ & 16 & Elements per dimension & --- \\
$N^2$ & 256 & Total elements per node & \textbf{new} \\
$R_0$ & $10^4$ & Nominal propagation distance ($\lambda$) & --- \\
$x_0$ & 2000 & Initial target $x$-coordinate & --- \\
$v$ & 1.0 & Target speed ($\lambda$/time) & --- \\
$\theta$ & 0 & Heading angle from $x$-axis & --- \\
$z_T$ & 10 & Constant $z$-offset (altitude) & --- \\
$M_z{\times}M_x$ & $20{\times}20$ & Pixel grid & --- \\
$2\delta$ & 10 & Pixel side ($\lambda$) & --- \\
$\xi_{\mathrm{th}}$ & 2.0 & Zone threshold & --- \\
SNR & 20\,dB & Per-element SNR & --- \\
\midrule
$\gamma$ & 5000 & Profile step size & --- \\
$\tau$ & $10^{-6}$ & TV regularisation & --- \\
$N_S$ & 200 & Profile iterations & --- \\
$N_D$ & 20 & Denoising iterations & --- \\
$N_V$ & 20 & Outer parameter iterations & --- \\
\bottomrule
\end{tabular}
\end{table}

%======================================================================
\section{Processing Pipeline Summary}
\label{sec:pipeline}
%======================================================================

\begin{enumerate}
  \item \textbf{FSSR Generation.}
    \begin{enumerate}
      \item For each time step~$t_m$: compute $R(t_m)=R_0-v_y t_m$
        and $\Delta x_k(t_m)=x_0+v_x t_m-\bar{X}_k$.
      \item For each pixel column~$q$: compute
        $R_{\mathrm{eff},q}(t_m)=R(t_m)-X'_q\tan\theta$.
      \item Compute $F_q^{(k)}(t_m)$, $G_{p,q}(t_m)$, and
        $H_{k,p,q}(t_m)$.
      \item True FSSR: $\varepsilon_k(t_m)=|1+\sum_{p,q}
        H_{k,p,q}P_{p,q}|^2$.
      \item Generate noisy signals~\eqref{eq:rx_signal} at all $N^2$
        elements; measure
        $b_k(t_m)=|\tilde{y}_k(t_m)|^2$~\eqref{eq:fssr_measured}
        via 2D beamforming.
    \end{enumerate}

  \item \textbf{Zone Classification.}
    $\xi_k(t)=|\Delta x_k(t)|/\sqrt{R(t)}$; Zone~B/A per node.

  \item \textbf{Zone~B.}
    \begin{enumerate}
      \item Pass~1: uncompensated covariance ($M\le
        M_{\max}^{(\mathrm{unc})}$); rough 2D DOA
        $(\hat{\ell}_x,\hat{\ell}_z)$ via separable or joint MUSIC.
      \item Pass~2+: phase-compensated covariance ($M\le
        M_{\mathrm{opt}}$); refined 2D DOA.
      \item Azimuth analysis: fit $\hat{\beta}_k$ vs.\ $\bar{X}_k$
        to extract $\hat{a},\hat{b}$.
      \item Doppler via STFT on 2D-steered beam; extract
        $\hat{\dot{f}}_k$.
      \item Estimate $\hat{v}_x,\hat{v}_y,\hat{R}_0,\hat{x}_0$.
      \item \textbf{Altitude from elevation DOA:}
        $\hat{z}_T^{(\mathrm{B})}=\hat{\alpha}_z\,\hat{R}_0$
        \hfill(\textbf{new in v6})
    \end{enumerate}

  \item \textbf{Zone~A.}
    \begin{enumerate}
      \item Collect $b_{k,\ell}$ via 2D beamforming at each node.
      \item Build coefficient matrices $\mathbf{H}_k(t_\ell)$ using
        $R_{\mathrm{eff},q}=\hat{R}(t_\ell)-X'_q\tan\hat{\theta}$.
      \item Algorithm~\ref{alg:joint_estimation}: joint retrieval of
        $\hat{\mathbf{P}}$ and refined parameters, \textbf{initialised
        with $\hat{z}_T^{(\mathrm{B})}$} from Zone~B elevation DOA.
    \end{enumerate}

  \item \textbf{Output.}
    $\hat{\mathbf{P}}\in\{0,1\}^{M_z\times M_x}$;\;
    $\hat{v}_x,\hat{v}_y,\hat{z}_T,\hat{R}_0$.
\end{enumerate}

%======================================================================
\section{Summary of Changes from v5.2}
\label{sec:changes}
%======================================================================

\begin{table}[H]
\centering
\caption{Summary of v5.2~$\to$~v6 changes}
\label{tab:changes}
\begin{tabular}{@{}p{3.5cm}p{5cm}p{5cm}@{}}
\toprule
\textbf{Component} & \textbf{v5.2} & \textbf{v6} \\
\midrule
Array geometry & $N$-element ULA along $x$ &
  $N\times N$ URA in $x$--$z$ plane \\
Elements per node & $N$ & $N^2$ \\
Beamformed noise & $\sigma^2/N$ & $\sigma^2/N^2$ \\
Measurable DOA & $\ell_x$ only &
  $(\ell_x,\ell_z)$ both \\
Zone~B $z_T$ estimation & $z_T=0$ default &
  $\hat{z}_T^{(\mathrm{B})}=\hat{\ell}_z\hat{R}_0$ from
  elevation DOA \\
Zone~A $z_T$ initialisation &
  $z_T^{(0)}=0$ &
  $z_T^{(0)}=\hat{z}_T^{(\mathrm{B})}$ \\
Steering vector & $\mathbf{a}_x(\ell_x)\in\C^N$ &
  $\mathbf{a}_x(\ell_x)\otimes\mathbf{a}_z(\ell_z)\in\C^{N^2}$ \\
DOA algorithms & 1D MUSIC, 1D LS-ESPRIT &
  2D MUSIC (joint/separable), 2D LS-ESPRIT \\
Steered beam (Doppler) & 1D steering &
  2D steering ($+10\log_{10}N$\,dB gain) \\
FSSR model & Unchanged & Unchanged \\
3D target model & Unchanged & Unchanged \\
Zone~A algorithms & Unchanged & Unchanged \\
\bottomrule
\end{tabular}
\end{table}

\begin{thebibliography}{9}

\bibitem{Paper1}
X.~Shen and D.~Huang,
``Pixel based shadow profile retrieval in forward scatter radar using
forward scatter shadow ratio,''
\emph{IEEE Access}, vol.~12, pp.~60467--60474, 2024.

\bibitem{Draft5}
X.~Shen and D.~Huang,
``A gradient-based method with denoising for target shadow profile
retrieval in forward scatter radar,''
submitted, 2025.

\bibitem{IETDraft4}
X.~Shen and D.~Huang,
``Retrieving target motion parameters together with target shadow
profile in forward scatter radar,''
submitted to \emph{IET Radar, Sonar \& Navigation}, 2025.

\end{thebibliography}

\end{document}
